\chapter{Despliegue, continuidad y seguridad}
\section{Ciclo de vida Docker}
El punto de entrada operativo es \code{scripts/stack.py}, invocado por Make. Todos los comandos comparten el proyecto original y los archivos Compose canónicos. El proyecto selecciona los volúmenes persistentes; cambiarlo para evitar un error puede conectar otra instalación y ocultar los registros antiguos.

\begin{longtable}{L{.23\linewidth}L{.67\linewidth}}
\caption{Comandos mantenidos y efecto sobre la instalación.}\label{tab:comandos}\\
\toprule Comando & Efecto \\ \midrule\endfirsthead
\toprule Comando & Efecto \\ \midrule\endhead
\code{make setup} / \code{make up} & Construyen lo necesario, aplican migraciones, comisionan o verifican el protocolo local y levantan la aplicación conservando identidades existentes.\\
\code{make simulate} & Levanta la aplicación completa y el simulador registrado; no requiere dejar una app parcial iniciada por otro archivo Compose.\\
\code{make doctor} / \code{make ps} & Inspeccionan dependencias, disco, despliegue, readiness y contenedores.\\
\code{make backup} & Exporta un archivo SQL privado y verifica su estructura con pg\_restore.\\
\code{make down} & Detiene servicios y conserva volúmenes.\\
\code{make recover} & Recupera una instalación local sobre el mismo PostgreSQL; si es necesario reemplaza cadena y archiva historia económica.\\
\code{make clean} & Elimina únicamente artefactos de compilación de contratos, no registros ni volúmenes.\\
\code{make demo-scenario} & Ejecuta o reanuda la prueba finita MQTT/API/EVM de 5/3 kWh.\\
\code{make blockchain-check} & Inserta telemetría validada directamente y verifica el mismo escenario; no prueba el transporte MQTT.\\
\bottomrule
\end{longtable}
Los servicios incluyen PostgreSQL/TimescaleDB, Redis, Mosquitto, Anvil, inicialización de estado, migración, API, worker, suscriptor, frontend y simulador opcional. API, worker y suscriptor esperan migración correcta y dependencias saludables. Los puertos públicos locales se limitan a loopback: 8080 para frontend, 8000 para API, 8545 para RPC y 1883 para MQTT.\cite{deployment,start}

\section{Persistencia EVM y recuperación}
La pérdida de Anvil fue una limitación operativa relevante: podían conservarse usuarios y lecturas mientras las direcciones almacenadas ya no contenían contratos. El error \code{BadFunctionCallOutput} y la ausencia de bytecode no se resuelven acreditando saldos en SQL.

El supervisor actual mantiene snapshots validados de hasta 128 MiB, con escritura en archivo temporal, fsync y reemplazo atómico; conserva un checkpoint anterior completo. Obtiene snapshots cada 30 segundos y durante la parada ordenada. Limita estados históricos y evita el crecimiento no controlado de logs mediante rotación. Si el snapshot actual falla, intenta el previo. Si ambos son inválidos, el arranque se detiene; no continúa silenciosamente en una cadena vacía.

Un corte abrupto puede perder operaciones posteriores al último checkpoint. Por eso conservar un archivo válido no elimina la necesidad de conciliación. La identidad de génesis se preserva y se comprueba; los relojes de minería se ajustan sobre el último estado cargado para no quedar atascados cerca de la fecha inicial.

Cuando la cadena es irrecuperable, \code{make recover}:
\begin{enumerate}
\item Detiene procesos que escriben y verifica acceso al mismo proyecto y base.
\item Crea un respaldo SQL privado verificable y conserva configuración y journals.
\item Archiva snapshots de Anvil y el historial financiero relacionado con la cadena perdida.
\item Registra una recuperación con referencias antiguas/nuevas y conserva las filas.
\item Despliega contratos locales nuevos y reafirma identidades y autorizaciones originales.
\item Fija un corte contable para evitar nueva emisión de energía histórica.
\item Conserva un marcador de progreso, checkpoint y configuración para poder reanudar una interrupción.
\end{enumerate}
Restaurar SQL no reconstruye almacenamiento EVM, recibos ni saldos. Los lotes archivados mantienen evidencia histórica; el usuario no debe esperar que una conciliación automática resucite una cadena borrada.

\section{Corrección de credenciales PostgreSQL}
La configuración previa forzaba \code{solar:solar} dentro del backend de Compose. Un volumen creado con otra contraseña podía aparecer saludable por pg\_isready, pero la migración fallaba con \code{InvalidPasswordError}. Se corrigió el despliegue para conservar la URL original y adaptar únicamente host y puerto a \code{timescaledb:5432} dentro del contenedor.

El módulo \code{src.core.docker_database} utiliza el parser de URL de SQLAlchemy, conserva credenciales codificadas, base y parámetros, y evita imprimir la URL en errores. El entrypoint exporta el resultado antes de la migración o la app. \code{POSTGRES_PASSWORD} inicializa volúmenes nuevos; editarlo no cambia la contraseña de un rol existente. La corrección de configuración no rota credenciales ni modifica registros de negocio.\cite{dbfix}

\section{Salud y diagnóstico}
\code{/health} comprueba vivacidad de API/base; \code{/ready} agrega Redis, red y código de registro, además de actividad de worker y suscriptor. Un estado ready indica disponibilidad actual, no que las lecturas sean válidas ni que un lote concreto se haya liquidado.

Las métricas del suscriptor incluyen persistidos, duplicados, rechazados y cola. Un máximo de hora de medición sin cambios entre dos consultas requiere compararlo con el reloj y los eventos recientes, no declarar inmediatamente detenido el simulador. Un mensaje rechazado puede ser antiguo y un contador de ventana móvil puede disminuir sin caída del proceso. El diagnóstico combina logs acotados, calidad por dispositivo, límites de intervalos y estados de lote/financiación.

\section{Seguridad implementada y límites}
Las contraseñas se almacenan con bcrypt; los JWT usan revocación de JTI en Redis y las verificaciones de wallet tienen nonce único. Las credenciales de medidor se comparan con digests y se excluyen del almacenamiento de evidencia. Los contenedores de backend operan sin root. Las firmas, nonces y génesis se registran antes de transmitir operaciones de servicio.

Persisten dependencias de confianza y necesidades de endurecimiento: tokens legibles por el navegador, refresh JWT sin rotación opaca, logout limitado al token presentado, correo de recuperación no desplegado como servicio canónico, gestión de secretos y alta disponibilidad pendientes. Los brokers anónimos locales y las claves de fixtures no son apropiados para un servicio público. Las operaciones sensibles de contratos usan permisos y protecciones de reentrancia, pero no existe una auditoría externa que certifique ausencia de vulnerabilidades.\cite{security}

\section{CI y reproducibilidad}
El workflow \code{.github/workflows/verify.yml} configura instalación de dependencias, suites de contratos/backend/IoT/frontend, chequeo de ABI, stack, escenario, reinicio del worker y pruebas de navegador. Su existencia demuestra la automatización configurada, no el éxito de una ejecución remota concreta; para afirmarlo se necesitaría el registro del job correspondiente.

Las dependencias se fijan con instaladores y lockfiles. La compilación y el despliegue emplean versiones distintas de herramientas de host y del nodo Docker según sus guías; no se supone que todos los ejecutables Foundry tengan la misma versión.
